% ===========================================================================
%  格点QCD GPU蒸馏计算管线 — 物理分析报告 (test0 / pyqcd)
%  Physical Analysis Report — test0 (pyqcd 调用)
% ===========================================================================
\documentclass[11pt,a4paper]{article}
\usepackage[UTF8]{ctex}
\setCJKmainfont{AR PL SungtiL GB}[BoldFont=AR PL UMing CN]
\setCJKsansfont{AR PL KaitiM GB}[BoldFont=AR PL UMing CN]
\setCJKmonofont{AR PL UMing CN}
\usepackage{amsmath,amssymb,amsfonts}
\usepackage{physics}
\usepackage{braket}
\usepackage{bm}
\usepackage[margin=2.5cm]{geometry}
\usepackage{booktabs}
\usepackage{graphicx}
\usepackage[colorlinks=true,linkcolor=blue,citecolor=blue,urlcolor=blue]{hyperref}
\usepackage{enumitem}
\usepackage{caption}
\usepackage{subcaption}
\usepackage{multirow}
\usepackage{array}

\newcommand{\gev}{\;\mathrm{GeV}}
\newcommand{\fm}{\;\mathrm{fm}}
\newcommand{\fmto}{\;\mathrm{fm}^{-1}}
\newcommand{\meff}{m_{\mathrm{eff}}}
\newcommand{\Nconf}{N_{\mathrm{conf}}}
\newcommand{\Nev}{N_{\mathrm{ev}}}
\newcommand{\Nt}{N_t}
\newcommand{\Nx}{N_x}
\newcommand{\tsep}{t_{\mathrm{sep}}}
\newcommand{\Ctwo}{C^{(2)}}
\newcommand{\Cthree}{C^{(3)}}
\newcommand{\Em}{E^{(0)}}
\newcommand{\pmom}{p_z}
\newcommand{\apm}{a^{-1}}
\newcommand{\jack}{\mathrm{JK}}
\newcommand{\gmu}{\gamma_\mu}

\title{\textbf{格点QCD GPU蒸馏计算管线物理分析报告}\\[0.3em]
       \large 顶点函数、Wick收缩、动态收缩与关联函数分析 (test0 — pyqcd)}
\author{张鑫\thanks{中国科学院近代物理研究所 (IMP, CAS)}}
\date{2026年08月15日}

\begin{document}
\maketitle

\begin{abstract}
本报告基于格点QCD蒸馏(Distillation)框架，在GPU (CUDA) 上实现了完整的关联函数计算管线：
顶点函数($VdV$/$VVV$)、Wick收缩分析、动态收缩、以及两点($pp$/$pn$)、OPE、三点($PJN$)、
四点($PJNNJNp$)关联函数，并进行Jackknife/有效质量/三点比值($ratio_{3p}$)统计分析。
本运行由 examples/test0 调用 pyqcd 包完成（与成功实例 docker-v20260805 逐项一致）。
计算使用CLQCD合作组的规范组态
($\beta=6.20$, $24^3\times72$, $a\approx0.1053\;\fm$, $\apm\approx1.874\;\gev$)，
共 1 个组态（6250），
计算精度 complex64。
\end{abstract}

\tableofcontents
\newpage

\section{引言}
LaMET (Large Momentum Effective Theory) 通过计算大动量下的准分布(quasi-distribution)
关联函数并做微扰匹配，得到光锥 parton 分布函数。胶子 PDF 涉及不相连(disconnected)图，
其中三点函数可分解为质子两点函数与胶子算符 (OPE) 两部分的乘积。本报告由 test0 套件
调用 pyqcd 包实现（逻辑照抄成功实例 docker-v20260805，自包含）。

\section{理论框架}
\subsection{格点系综参数}
\begin{table}[h]
\centering
\caption{格点系综参数 (表~\ref{tab:ensemble})}
\label{tab:ensemble}
\begin{tabular}{ll}
\toprule
参数 & 值 \\
\midrule
$\beta$ & 6.20 (Clover Wilson) \\
格点 & $24^3\times72$ \\
格距 $a$ & 0.1053 fm \\
逆格距 $\apm$ & $\approx 1.874$ GeV \\
本征矢数 $\Nev$ / $N_{\mathrm{ev,1}}$ & 100 / 100 \\
动量 & $P=(0,0,0)$, $P=(0,0,2)$ \\
组态数 $\Nconf$ & 1 \\
精度 & complex64 \\
\bottomrule
\end{tabular}
\end{table}

\subsection{蒸馏方法与顶点函数}
蒸馏 (distillation) 方法把夸克传播子投影到拉普拉斯算符的低模空间：
\[ \tau_{ij}(t_s,t_f) = v^\dagger_i(t_s)\, M^{-1}(t_s,t_f)\, v_j(t_f). \]
两点关联函数所需的顶点函数为
\begin{equation}
V^{VdV}_{mn}(\mathbf{p}) = \sum_{\mathbf{x}} e^{-i\mathbf{p}\cdot\mathbf{x}}\,
    v^\dagger_m(\mathbf{x})\, v_n(\mathbf{x}),
\label{eq:VdV}
\end{equation}
\begin{equation}
V^{VVV}_{m n l}(\mathbf{p}) = \sum_{\mathbf{x}} e^{-i\mathbf{p}\cdot\mathbf{x}}\,
    \varepsilon_{abc}\, v^a_m(\mathbf{x})\, v^b_n(\mathbf{x})\, v^c_l(\mathbf{x}),
\label{eq:VVV}
\end{equation}
其中 $VdV$ 用于介子，$VVV$ 用于重子（质子/中子）。

\subsection{两点关联函数与有效质量}
源平均后的两点函数为
\[ C(t) = \frac{1}{\Nt}\sum_{t_s} C(t_s,\, t_s + t). \]
有效质量的对数形式与双曲余弦形式为
\begin{equation}
\meff(t) = \ln\frac{C(t)}{C(t+1)}\cdot\frac{\hbar c}{a}, \qquad
\meff(t) = \operatorname{arccosh}\frac{C(t+2)+C(t)}{2C(t+1)}\cdot\frac{\hbar c}{a}.
\label{eq:meff}
\end{equation}

\subsection{三点/两点比值}
连通的三点函数 $C^{(3)}(\tau)$（质子-矢量流-核子）与两点函数 $C^{(2)}$ 的比值采用
lqcddb 的 $ratio_{3p}$ 公式，包含 $\sqrt{\cdots}$ 因子。不相连胶子比值则按
huangcl 的 code\_1.py 算法构造：
\begin{equation}
C^{(3)}(t, \tau, z) = C^{(2)}(t)\, O(z,\tau),
\qquad R(t,\tau,z) = \frac{C^{(3)} - C^{(2)}\braket{O}}{C^{(2)}},
\end{equation}
并对每个 $z$ 做关联拟合 $R(t,\tau) = c_0 + c_1 e^{-dE\,\tau} + c_1 e^{-dE\,(t-\tau)}$。

\section{计算方法 (GPU管线)}
管线步骤：
\begin{enumerate}
    \item 顶点函数：$VdV$/$VVV$（GPU，按时间片流式计算）
    \item Wick收缩分析 + 动态收缩（lqcddb 引擎，注册表 + einsum 计划缓存）
    \item 关联函数：2pt ($pp$/$pn$/pion)、OPE (胶子算符)、3pt ($PJN$)、4pt ($PJNNJNp$)
    \item 统计分析：Jackknife、有效质量、$ratio_{3p}$（code\_1.py 形式）
    \item 绘图与 LaTeX 报告
\end{enumerate}
全部中间结果与日志均保存在版本目录中。

\section{结果与分析}
\subsection{两点关联函数与有效质量}
表~\ref{tab:meff} 给出各道的有效质量（Jackknife，加权平台）。
\begin{table}[h]
\centering
\caption{有效质量 (加权平台) (表~\ref{tab:meff})}
\label{tab:meff}
\begin{tabular}{llccccl}
\toprule
粒子 & 动量 & $E_0$ [GeV] & 期望 [GeV] & 平台 & 点数 \\
\midrule
    proton & $P={P0}$ & — & 1.000 & $[2,8]$ & 0.0 \\
    proton & $P={P2}$ & — & — & $[2,8]$ & 0.0 \\
    pion & $P={P0}$ & — & 0.300 & $[2,8]$ & 0.0 \\
    pion & $P={P2}$ & — & — & $[2,8]$ & 0.0 \\
\bottomrule
\end{tabular}
\end{table}

\subsection{三点/两点连通比值}
表~\ref{tab:ratio} 给出 $\gamma_3$（$z$方向）分量的比值 $R(\tau)$。
\begin{table}[h]
\centering
\caption{连通三点/两点比值 (表~\ref{tab:ratio})}
\label{tab:ratio}
\begin{tabular}{llc}
\toprule
粒子 & 动量 & $R(\tau{\approx}4)$ \\
\midrule
    proton & $P={P0}$ & $+nan \pm nan$ \\
    proton & $P={P2}$ & $+nan \pm nan$ \\
    pion & $P={P0}$ & $+nan \pm nan$ \\
    pion & $P={P2}$ & $+nan \pm nan$ \\
\bottomrule
\end{tabular}
\end{table}

\subsection{不相连胶子比值与拟合}
表~\ref{tab:disc} 给出质子 $P_z=2$ 道不相连比值按 code\_1.py 拟合的参数。
\begin{table}[h]
\centering
\caption{不相连比值拟合参数 $R=c_0+c_1e^{-dE\,\tau}+c_1e^{-dE\,(t-\tau)}$ (表~\ref{tab:disc})}
\label{tab:disc}
\begin{tabular}{lcccc}
\toprule
$z$ & $c_0$ & $c_1$ & $dE$ & $\chi^2/\mathrm{dof}$ \\
\midrule

\bottomrule
\end{tabular}
\end{table}

\subsection{组态一致性}
表~\ref{tab:config} 给出各道 $C(0)$ 的组态间离散程度。
\begin{table}[h]
\centering
\caption{各道 $C(0)$ 组态一致性 (表~\ref{tab:config})}
\label{tab:config}
\begin{tabular}{lcc}
\toprule
道 & $\braket{C(0)}$ & 相对离散度 \\
\midrule

\bottomrule
\end{tabular}
\end{table}

\subsection{计算耗时}
\begin{table}[h]
\centering
\caption{分步耗时 (表~\ref{tab:timing})}
\label{tab:timing}
\begin{tabular}{lr}
\toprule
步骤 & 耗时 \\
\midrule
    analysis & 0.2 s \\
    plots & 1.7 s \\
\bottomrule
\end{tabular}
\end{table}

\section{讨论与展望}
当前运行使用单精度 (complex64)、$\Nev=100$、$\Nconf=1$。
$pn$（质子-中子）两点函数因味守恒而恒为零（质子 $uud$ 与中子 $udd$ 味结构不同），
这与理论预期一致。动量 $P=(0,0,2)$ 对应物理动量
$p_z = \frac{2\pi\cdot 2}{24\,a}\approx 0.981\;\gev$。后续工作可增加本征矢数目、
组态数目、动量涂抹、多源与 GEVP 以改善激发态污染。

\section{结论}
\begin{enumerate}
    \item 完整实现了从顶点函数到四点关联函数的 GPU 蒸馏管线（pyqcd 调用）。
    \item 两点函数与有效质量分析通过 Jackknife 获得统计误差。
    \item 三点 ($PJN$) 与四点 ($PJNNJNp$) 关联函数及比值分析完成。
    \item OPE 胶子算符按 donghx 算法计算并与两点函数组合成不相连比值。
\end{enumerate}

\begin{thebibliography}{9}
\bibitem{zhang2019} J.-H. Zhang et al., PRL 122, 142001 (2019).
\bibitem{fan2021} Z. Fan et al., PRD 104, 074502 (2021).
\bibitem{ji2013} X. Ji, PRL 110, 262002 (2013).
\bibitem{peardon2009} M. Peardon et al., PRD 80, 054506 (2009).
\end{thebibliography}

\end{document}
